% ===========================================================
% Chapter 16: Multi-Agent Systems and RL
% Part IV: Frontiers
% ===========================================================

\chapter{Multi-Agent Systems and RL}
\label{ch:agents}

\begin{quote}
\itshape
The whole is more than the sum of its parts.
\par\medskip
\upshape\raggedleft
---\,Aristotle, \textit{Metaphysics}
\end{quote}

\bigskip

Every setting we have studied so far involves a single learner interacting
with a world whose dynamics are fixed. The environment may be
stochastic, partially observed, or high-dimensional, but it does not
itself \emph{learn}. In the real world, this assumption rarely holds.
Financial markets are populated by traders who each adapt to the
strategies of the others. Autonomous vehicles must negotiate
intersections with other vehicles that are simultaneously trying to
navigate safely. A team of rescue robots must divide a building among
themselves without central coordination. Language models deployed as
assistants increasingly delegate sub-tasks to specialised
agents\,---\,other models that respond, check, or criticise. In every
case, the ``environment'' contains \emph{other agents}, and those agents
are themselves learning and adapting.

Multi-agent reinforcement learning (MARL)\index{multi-agent reinforcement learning|textbf}
is the study of how multiple decision-making entities can learn good
behaviour when they coexist and interact. The field sits at the
intersection of game theory, optimisation, and modern deep learning,
and it poses challenges that are qualitatively different from
single-agent RL: the environment is \emph{non-stationary} because
other agents are changing; rewards may be shared, competing, or
somewhere in between; and the space of joint behaviours is
exponentially larger than the space of individual behaviours.

This chapter develops the theory and practice of MARL progressively.
Section~\ref{sec:markov-games} introduces Markov games\,---\,the
natural multi-agent generalisation of MDPs\,---\,and defines the
solution concept of Nash equilibrium. Section~\ref{sec:cooperative}
covers cooperative MARL, including the centralised-training
with decentralised-execution (CTDE) paradigm and the QMIX and MAPPO
algorithms. Section~\ref{sec:competitive} studies competitive and
mixed settings, including zero-sum games, self-play, and the story of
AlphaGo and AlphaZero. Section~\ref{sec:communication} examines
emergent communication and social dilemmas.
Section~\ref{sec:marl-llms} brings the framework to large language
models, covering debate, multi-agent code generation, and agentic
workflows. Section~\ref{sec:marl-challenges} catalogues the open
challenges, and Section~\ref{sec:marl-summary} summarises.

% ===========================================================
\section{Multi-Agent Markov Games}
\label{sec:markov-games}
% ===========================================================

\subsection{From MDPs to Markov Games}

Recall from Chapter~\ref{ch:mdp} that a Markov decision process is
a tuple $(\cS, \cA, \cT, \cR, \gamma)$ in which a single agent takes
actions, the environment transitions to a new state, and the agent
receives a scalar reward. When there are $N \geq 2$ agents, the
natural extension is a \emph{Markov game}\index{Markov game|textbf},
also called a \emph{stochastic game}\index{stochastic game}.

\begin{definition}[Markov Game]
\label{def:markov-game}
A \textbf{Markov game} with $N$ agents is a tuple
\[
  \mathcal{G} = \bigl(\cS,\, \cA^1, \ldots, \cA^N,\, \cT,\,
                       R^1, \ldots, R^N,\, \gamma\bigr),
\]
where:
\begin{itemize}[leftmargin=*]
  \item $\cS$ is a finite set of \textbf{states} shared by all agents.
  \item $\cA^i$ is the \textbf{action set} for agent $i$. The
    \textbf{joint action space} is $\cA = \cA^1 \times \cdots \times \cA^N$.
  \item $\cT : \cS \times \cA \to \Delta(\cS)$ is the
    \textbf{transition kernel}; the next state depends on the current
    state \emph{and} the joint action of all agents.
  \item $R^i : \cS \times \cA \times \cS \to \R$ is the
    \textbf{reward function} for agent $i$. Each agent has its own
    reward, which may depend on the actions of all agents.
  \item $\gamma \in [0,1)$ is a shared discount factor.
\end{itemize}
\end{definition}

When $N = 1$, a Markov game reduces to a standard MDP. When all
agents share the same reward function, $R^1 = \cdots = R^N$, the game
is \textbf{fully cooperative}\index{fully cooperative game}. When
$\sum_i R^i = 0$ at every transition, it is \textbf{zero-sum}\index{zero-sum game}.
The general case, where rewards are neither identical nor summing to
zero, is called a \textbf{mixed or general-sum
game}\index{general-sum game}.

The joint policy\index{joint policy} of all agents is
$\boldsymbol{\pi} = (\pi^1, \ldots, \pi^N)$, where each $\pi^i :
\cS \to \Delta(\cA^i)$ is a stochastic policy for agent $i$.
Agent $i$'s objective is to maximise its own discounted return:
\begin{equation}
\label{eq:marl-return}
J^i(\boldsymbol{\pi}) = \E_{\boldsymbol{\pi}}\!\left[
  \sum_{t=0}^{\infty} \gamma^t R^i(s_t, \bm{a}_t, s_{t+1})
\right],
\end{equation}
where $\bm{a}_t = (a_t^1, \ldots, a_t^N)$ is drawn from the product
policy $\prod_i \pi^i(\cdot \mid s_t)$.



\begin{figure}[ht]
\centering
\begin{tikzpicture}[
  state/.style={circle, draw, minimum size=1.0cm, font=\small},
  agent/.style={rectangle, draw, rounded corners, minimum width=1.8cm,
                minimum height=0.7cm, font=\small, fill=gray!12},
  action/.style={->, >=Stealth, thick},
  reward/.style={->, >=Stealth, dashed, red!70!black},
  trans/.style={->, >=Stealth, thick, blue!60!black}
]

% States
\node[state] (s) at (0,0) {$s_t$};
\node[state] (sn) at (5,0) {$s_{t+1}$};

% Agents
\node[agent] (a1) at (-2.2, 1.5) {Agent $1$};
\node[agent] (aN) at (-2.2,-1.5) {Agent $N$};
\node[font=\small] (dots) at (-2.2, 0) {$\vdots$};

% Actions into state
\draw[action] (a1.east) -- node[above, font=\scriptsize] {$a_t^1$}
              ([yshift=3pt]s.west);
\draw[action] (aN.east) -- node[below, font=\scriptsize] {$a_t^N$}
              ([yshift=-3pt]s.west);

% Transition
\draw[trans] (s) -- node[above, font=\small] {$\mathcal{T}(s_{t+1}\mid s_t,\bm{a}_t)$} (sn);

% Reward arrows: orthogonal with one right angle
\draw[reward] (sn.north) -- ++(0, 1.2) -| (a1.east)
              node[pos=0.25, above, font=\scriptsize, red!70!black] {$R^1$};
\draw[reward] (sn.south) -- ++(0,-1.2) -| (aN.east)
              node[pos=0.25, below, font=\scriptsize, red!70!black] {$R^N$};

\end{tikzpicture}
\caption{A Markov game with $N$ agents. At each timestep, every agent
  observes the shared state $s_t$ and chooses an action. The joint
  action $\bm{a}_t$ drives the transition to $s_{t+1}$, and each
  agent receives its own individual reward $R^i$.}
\label{fig:markov-game}
\end{figure}
% \begin{figure}[ht]
% \centering
% \begin{tikzpicture}[
%   node distance=1.6cm and 2.4cm,
%   state/.style={circle, draw, minimum size=1.0cm, font=\small},
%   agent/.style={rectangle, draw, rounded corners, minimum width=1.8cm,
%                 minimum height=0.7cm, font=\small, fill=gray!12},
%   action/.style={->, >=Stealth, thick},
%   reward/.style={->, >=Stealth, dashed, red!70!black},
%   trans/.style={->, >=Stealth, thick, blue!60!black}
% ]
% % State node
% \node[state] (s) {$s_t$};
% \node[state, right=3.8cm of s] (sn) {$s_{t+1}$};

% % Agents
% \node[agent, above left=0.5cm and 0.4cm of s] (a1) {Agent $1$};
% \node[agent, below left=0.5cm and 0.4cm of s] (aN) {Agent $N$};
% \node[font=\small, left=0.55cm of s, yshift=0.05cm] (dots) {$\vdots$};

% % Actions into state
% \draw[action] (a1.east) -- node[above, font=\scriptsize] {$a_t^1$}
%               ([yshift=3pt]s.west);
% \draw[action] (aN.east) -- node[below, font=\scriptsize] {$a_t^N$}
%               ([yshift=-3pt]s.west);

% % Transition
% \draw[trans] (s) -- node[above, font=\small] {$\cT(s_{t+1}\mid s_t,\bm{a}_t)$} (sn);

% % Reward arrows
% \draw[reward] (sn.north) to[bend left=30]
%               node[above right, font=\scriptsize, red!70!black] {$R^1$} (a1.east);
% \draw[reward] (sn.south) to[bend right=30]
%               node[below right, font=\scriptsize, red!70!black] {$R^N$} (aN.east);

% \end{tikzpicture}
% \caption{A Markov game with $N$ agents. At each timestep, every agent
%   observes the shared state $s_t$ and chooses an action. The joint
%   action $\bm{a}_t$ drives the transition to $s_{t+1}$, and each
%   agent receives its own individual reward $R^i$.}
% \label{fig:markov-game}
% \end{figure}

\subsection{Nash Equilibrium}

Because each agent's return depends on the policies of \emph{all}
agents, the standard notion of optimality from single-agent RL no
longer applies. Raising agent $i$'s return by changing $\pi^i$ may
require the other agents to hold their policies fixed\,---\,or it may
backfire if they adapt. The solution concept borrowed from game theory
is the \textbf{Nash equilibrium}\index{Nash equilibrium|textbf}.

\begin{definition}[Nash Equilibrium]
\label{def:nash}
A joint policy $\boldsymbol{\pi}^* = (\pi^{1*}, \ldots, \pi^{N*})$
is a \textbf{Nash equilibrium} if, for every agent $i$ and every
alternative policy $\tilde\pi^i$,
\[
  J^i(\pi^{i*}, \boldsymbol{\pi}^{-i*}) \;\geq\;
  J^i(\tilde\pi^i, \boldsymbol{\pi}^{-i*}),
\]
where $\boldsymbol{\pi}^{-i*}$ denotes the policies of all agents
except $i$, held fixed at their equilibrium values.
\end{definition}

In words, a Nash equilibrium is a joint policy from which no single
agent has an incentive to deviate unilaterally. Nash's theorem
guarantees that every finite game has at least one Nash equilibrium
in mixed strategies, but equilibria need not be unique, and finding
one is computationally hard in general ($\mathsf{PPAD}$-complete for
two-player general-sum games).

\begin{example}[Prisoner's Dilemma as a One-Step Markov Game]
\label{ex:prisoners}
Two agents simultaneously choose \textsc{Cooperate} ($C$) or
\textsc{Defect} ($D$). Payoffs (row = Agent 1, column = Agent 2):
\begin{center}
\small
\begin{tabular}{@{}lcc@{}}
\toprule
 & $C$ & $D$ \\
\midrule
$C$ & $(3,3)$ & $(0,5)$ \\
$D$ & $(5,0)$ & $(1,1)$ \\
\bottomrule
\end{tabular}
\end{center}
Agent 1 always prefers $D$ regardless of Agent 2's action: if Agent 2
plays $C$, defecting yields $5 > 3$; if Agent 2 plays $D$, defecting
yields $1 > 0$. The same logic applies to Agent 2. Therefore $(D, D)$
is the unique Nash equilibrium, even though $(C, C)$ would give both
agents a higher payoff of $3$. This tension between individual
rationality and collective welfare is the hallmark of social
dilemmas\index{social dilemma} and will reappear in
Section~\ref{sec:communication}.
\end{example}

\subsection{Partial Observability and the Dec-POMDP}

In practice, agents rarely observe the full global state. Each agent
$i$ instead receives a local observation $o^i_t$ drawn from an
observation function $O^i : \cS \to \Delta(\mathcal{O}^i)$. The
resulting framework is a \textbf{decentralised partially observable
MDP}\index{Dec-POMDP} (Dec-POMDP), which is $\mathsf{NEXP}$-hard to
solve in the worst case. Much of practical MARL research focuses on
restricting the setting\,---\,cooperative teams, shared observations
at training time, or limited communication\,---\,to make the problem
tractable.

% ===========================================================
\section{Cooperative MARL}
\label{sec:cooperative}
% ===========================================================

\subsection{The Cooperative Setting}

In fully cooperative MARL all agents share a single team reward
$r_t = R(s_t, \bm{a}_t, s_{t+1})$. The joint objective is simply
\begin{equation}
\label{eq:cooperative-return}
J(\boldsymbol{\pi}) = \E_{\boldsymbol{\pi}}\!\left[
  \sum_{t=0}^{\infty} \gamma^t r_t \right].
\end{equation}
Despite the common objective, each agent must act on its own
local observation and cannot, at execution time, communicate or
observe the global state. This is the standard setting in
multi-robot coordination, multi-player cooperative games (e.g.\
StarCraft~II micromanagement), and multi-agent dialogue systems.

\subsection{Centralised Training with Decentralised Execution}

The most influential paradigm for cooperative MARL is
\textbf{centralised training with decentralised
execution}\index{centralised training with decentralised execution|textbf}
(CTDE)\index{CTDE|see{centralised training with decentralised execution}}.
During \emph{training}, a central critic or mixer has access to the
global state $s_t$ and the actions of all agents. During
\emph{execution}, each agent relies only on its local observation
history $\tau^i = (o^i_0, a^i_0, \ldots, o^i_t)$. The training
privilege allows the critic to give informative gradient signal; the
execution constraint ensures the policy is deployable in settings
where a central controller is unavailable.


\begin{figure}[ht]
\centering
\begin{tikzpicture}[
  box/.style={rectangle, draw, rounded corners, minimum width=2.0cm,
              minimum height=0.6cm, font=\small, fill=gray!10, align=center},
  mixer/.style={rectangle, draw, rounded corners, minimum width=2.4cm,
               minimum height=0.7cm, font=\small, fill=blue!12, align=center},
  envbox/.style={rectangle, draw, rounded corners, minimum width=2.4cm,
               minimum height=0.7cm, font=\small, fill=green!8, align=center},
  lossnode/.style={font=\small},
  phase/.style={rectangle, draw, dashed, rounded corners,
                inner sep=10pt},
  arr/.style={->, >=Stealth, thick},
  darr/.style={->, >=Stealth, dashed, gray}
]

%% ============ TRAINING PHASE ============

\node[box] (q1) at (0,0) {$Q^1(\tau^1, a^1)$};
\node[box] (q2) at (0,-1.2) {$Q^2(\tau^2, a^2)$};
\node[font=\small] (qdots) at (0,-2.1) {$\vdots$};
\node[box] (qN) at (0,-2.8) {$Q^N(\tau^N, a^N)$};

\node[mixer] (mix) at (5.0,-1.2) {\textbf{Mixer} $Q_\mathrm{tot}(s,\bm{a})$};
\node[font=\small] (glob) at (5.0,0.2) {global state $s$};
\node[lossnode] (loss) at (7.8,-1.2) {$\mathcal{L}_\mathrm{TD}$};

% Q^i -> Mixer: use explicit coordinates to guarantee entry into mixer
\draw[arr] (q1.east) -- (mix.west |- q1.east) -- (mix.west);
\draw[arr] (q2.east) -- (mix.west);
\draw[arr] (qN.east) -- (mix.west |- qN.east) -- (mix.west);

% Global state -> Mixer
\draw[arr, dashed] (glob.south) -- (mix.north);

% Mixer -> Loss
\draw[arr] (mix.east) -- (loss.west);

\node[phase, fit=(q1)(qN)(glob)(loss),
      label={[font=\small\bfseries]above:Training (centralised)}] (trainbox) {};

%% ============ EXECUTION PHASE ============

\node[box] (e1) at (0,-5.2) {Policy $\pi^1(\tau^1)$};
\node[box] (e2) at (0,-6.4) {Policy $\pi^2(\tau^2)$};
\node[font=\small] (edots) at (0,-7.3) {$\vdots$};
\node[box] (eN) at (0,-8.0) {Policy $\pi^N(\tau^N)$};

\node[envbox] (env) at (5.0,-6.4) {Environment};

% Policies -> Environment: same pattern
\draw[arr] (e1.east) -- (env.west |- e1.east) -- (env.west);
\draw[arr] (e2.east) -- (env.west);
\draw[arr] (eN.east) -- (env.west |- eN.east) -- (env.west);

\node[phase, fit=(e1)(eN)(env),
      label={[font=\small\bfseries]above:Execution (decentralised)}] (execbox) {};

%% ============ DASHED LINKS: Q-nets -> Policies ============
\draw[darr] (q1.south) -- (e1.north);
\draw[darr] (q2.south) -- (e2.north);
\draw[darr] (qN.south) -- (eN.north);

%% ============ VERTICAL DOTTED SEPARATOR ============
\coordinate (septop) at (2.5, 0.9);
\coordinate (sepbot) at (2.5,-8.7);
\draw[dotted, thick, gray!60] (septop) -- (sepbot);

\end{tikzpicture}
\caption{The CTDE architecture. During training a \emph{mixer}
  (e.g.\ QMIX) combines per-agent utility networks $Q^i$ into a
  joint $Q_\mathrm{tot}$ using the global state, and gradients flow
  back through all $Q^i$. At execution time the centralised mixer is
  discarded; each agent acts greedily with respect to its own $Q^i$
  using only its local observation history $\tau^i$.}
\label{fig:ctde}
\end{figure}
% \begin{figure}[ht]
% \centering
% \begin{tikzpicture}[
%   box/.style={rectangle, draw, rounded corners, minimum width=2.0cm,
%               minimum height=0.6cm, font=\small, fill=gray!10, align=center},
%   mixer/.style={rectangle, draw, rounded corners, minimum width=2.4cm,
%                minimum height=0.7cm, font=\small, fill=blue!12, align=center},
%   envbox/.style={rectangle, draw, rounded corners, minimum width=2.4cm,
%                minimum height=0.7cm, font=\small, fill=green!8, align=center},
%   lossnode/.style={font=\small},
%   phase/.style={rectangle, draw, dashed, rounded corners,
%                 inner sep=10pt},
%   arr/.style={->, >=Stealth, thick},
%   darr/.style={->, >=Stealth, dashed, gray}
% ]

% %% ============ TRAINING PHASE (top half) ============

% % Per-agent Q-value networks (left column)
% \node[box] (q1) at (0,0) {$Q^1(\tau^1, a^1)$};
% \node[box] (q2) at (0,-1.0) {$Q^2(\tau^2, a^2)$};
% \node[font=\small] (qdots) at (0,-1.7) {$\vdots$};
% \node[box] (qN) at (0,-2.3) {$Q^N(\tau^N, a^N)$};

% % Mixer (right of Q-column, vertically centred)
% \node[mixer] (mix) at (4.5,-1.0) {\textbf{Mixer} $Q_\mathrm{tot}(s,\bm{a})$};
% \node[font=\small] (glob) at (4.5,0) {global state $s$};

% % Loss (right of mixer)
% \node[lossnode] (loss) at (7.2,-1.0) {$\mathcal{L}_\mathrm{TD}$};

% % Arrows: Q^i -> Mixer (horizontal, no crossings)
% \draw[arr] (q1.east) -- (q1.east -| mix.west);
% \draw[arr] (q2.east) -- (mix.west);
% \draw[arr] (qN.east) -- (qN.east -| mix.west);
% \draw[arr, dashed] (glob) -- (mix);

% % Mixer -> Loss
% \draw[arr] (mix.east) -- (loss.west);

% % Training bounding box
% \node[phase, fit=(q1)(qN)(glob)(loss),
%       label={[font=\small\bfseries]above:Training (centralised)}] (trainbox) {};

% %% ============ EXECUTION PHASE (bottom half) ============

% % Per-agent policies (left column, well below training)
% \node[box] (e1) at (0,-4.8) {Policy $\pi^1(\tau^1)$};
% \node[box] (e2) at (0,-5.8) {Policy $\pi^2(\tau^2)$};
% \node[font=\small] (edots) at (0,-6.5) {$\vdots$};
% \node[box] (eN) at (0,-7.1) {Policy $\pi^N(\tau^N)$};

% % Environment (right of policy column, vertically centred)
% \node[envbox] (env) at (4.5,-5.8) {Environment};

% % Arrows: policies -> environment (horizontal, no crossings)
% \draw[arr] (e1.east) -- (e1.east -| env.west);
% \draw[arr] (e2.east) -- (env.west);
% \draw[arr] (eN.east) -- (eN.east -| env.west);

% % Execution bounding box
% \node[phase, fit=(e1)(eN)(env),
%       label={[font=\small\bfseries]above:Execution (decentralised)}] (execbox) {};

% %% ============ DASHED LINKS: Q-nets -> Policies ============
% \draw[darr] (q1.south) -- (e1.north);
% \draw[darr] (q2.south) -- (e2.north);
% \draw[darr] (qN.south) -- (eN.north);

% %% ============ VERTICAL DOTTED SEPARATOR ============
% \coordinate (septop) at ($(q1.north east)!0.5!(q1.north east -| mix.north west) + (0, 0.6)$);
% \coordinate (sepbot) at ($(eN.south east)!0.5!(eN.south east -| env.south west) + (0, -0.6)$);
% \draw[dotted, thick, gray!60] (septop) -- (sepbot);

% \end{tikzpicture}
% \caption{The CTDE architecture. During training a \emph{mixer}
%   (e.g.\ QMIX) combines per-agent utility networks $Q^i$ into a
%   joint $Q_\mathrm{tot}$ using the global state, and gradients flow
%   back through all $Q^i$. At execution time the centralised mixer is
%   discarded; each agent acts greedily with respect to its own $Q^i$
%   using only its local observation history $\tau^i$.}
% \label{fig:ctde}
% \end{figure}

\subsection{QMIX}

\textbf{QMIX}\index{QMIX|textbf} (Rashid et al., 2018) is the
canonical CTDE algorithm for cooperative value-based MARL. Each
agent $i$ maintains a per-agent utility network
$Q^i_\theta(\tau^i, a^i)$ that conditions only on the local
observation history. A central \emph{mixer network} combines these
utilities into a joint action-value function:
\begin{equation}
\label{eq:qmix}
Q_\mathrm{tot}(s, \bm{a}) = f_\phi\!\bigl(Q^1(\tau^1, a^1), \ldots,
  Q^N(\tau^N, a^N),\; s\bigr).
\end{equation}
The mixer uses hypernetworks\index{hypernetwork} conditioned on the
global state $s$ to produce non-negative weights, which ensures
the \textbf{monotonicity constraint}\index{QMIX!monotonicity}:
\[
  \frac{\partial Q_\mathrm{tot}}{\partial Q^i} \geq 0
  \quad \text{for all } i.
\]
This constraint is the key insight. It guarantees the
\textbf{Individual-Global-Max}\index{IGM condition} (IGM) property:
\[
  \argmax_{\bm{a}}\, Q_\mathrm{tot}(s, \bm{a})
  \;=\; \bigl(\argmax_{a^1} Q^1(\tau^1, a^1),\; \ldots,\;
               \argmax_{a^N} Q^N(\tau^N, a^N)\bigr).
\]
Maximising $Q_\mathrm{tot}$ is therefore equivalent to each agent
greedily choosing its own best action\,---\,no joint-action search
is needed at execution time. The whole system is trained end-to-end
with a standard TD loss:
\begin{equation}
\label{eq:qmix-loss}
\mathcal{L}(\theta, \phi) = \E\!\left[
  \bigl(r + \gamma\, \max_{\bm{a}'} Q_\mathrm{tot}(s', \bm{a}')
  - Q_\mathrm{tot}(s, \bm{a})\bigr)^2
\right].
\end{equation}
QMIX achieves state-of-the-art performance on the StarCraft
Multi-Agent Challenge (SMAC)\index{SMAC} benchmark, a suite of
cooperative micromanagement tasks in StarCraft~II.

\subsection{MAPPO}

While QMIX operates in the value-function paradigm, many cooperative
tasks benefit from policy-gradient methods because they handle
continuous actions naturally and are easier to tune. \textbf{MAPPO}
\index{MAPPO|textbf} (Yu et al., 2022) applies Proximal Policy
Optimisation (PPO; Chapter~\ref{ch:pg}) in the CTDE framework.
Each agent $i$ maintains a separate actor $\pi^i_\theta(\cdot \mid
\tau^i)$ conditioned on its local observation history. A \emph{shared
centralised critic} $V_\phi(s)$ conditions on the global state and
produces a baseline for variance reduction:
\begin{equation}
\label{eq:mappo-objective}
\mathcal{L}_\mathrm{MAPPO}(\theta^i) =
\E\!\left[\min\!\bigl(r^i_t(\theta^i)\,\hat{A}_t,\;
  \mathrm{clip}(r^i_t(\theta^i), 1-\epsilon, 1+\epsilon)\,\hat{A}_t
\bigr)\right],
\end{equation}
where $r^i_t(\theta^i) = \pi^i_\theta(a^i_t \mid \tau^i_t) /
\pi^i_{\mathrm{old}}(a^i_t \mid \tau^i_t)$ is the probability ratio
and $\hat{A}_t$ is the advantage estimated by the centralised critic.
A key practical finding is that \emph{parameter sharing}\index{parameter
sharing}\,---\,initialising all agents from the same actor network
while conditioning on an agent-index feature\,---\,drastically
improves sample efficiency by letting agents share gradients.

\begin{remark}[Parameter Sharing vs.\ Heterogeneous Agents]
Parameter sharing is powerful when agents play symmetric roles, but
can be inappropriate for teams of heterogeneous specialists (e.g.\ a
healer and a tank in a game). In that setting, separate networks with
a shared feature extractor often work best.
\end{remark}

% ===========================================================
\section{Competitive and Mixed Settings}
\label{sec:competitive}
% ===========================================================

\subsection{Zero-Sum Games and Self-Play}

In a two-player zero-sum game, $R^1 = -R^2$, so one agent's gain is
exactly the other's loss. The Nash equilibrium of a zero-sum Markov
game corresponds to the minimax policy: agent 1 maximises while agent
2 minimises the same cumulative reward. Formally, the minimax value
$v^*$ satisfies
\[
  v^* = \max_{\pi^1} \min_{\pi^2}\; J^1(\pi^1, \pi^2)
       = \min_{\pi^2} \max_{\pi^1}\; J^1(\pi^1, \pi^2),
\]
where the equality (von Neumann's minimax theorem)\index{minimax theorem}
holds when policies are allowed to be mixed (stochastic).

Training an agent to play optimally in a zero-sum game faces an
immediate bootstrapping problem: to learn to defeat strong opponents,
the agent needs strong opponents to learn against. \textbf{Self-play}
\index{self-play|textbf} resolves this by having the agent play against
past versions of itself. A canonical training loop is:
\begin{enumerate}[leftmargin=*]
  \item Maintain a population (or a single snapshot) of past policy
    checkpoints.
  \item Sample an opponent from the population.
  \item Collect rollouts; update the current policy against this
    opponent.
  \item Add the current policy to the population periodically.
\end{enumerate}
Self-play guarantees that the opponent distribution is always at a
difficulty level commensurate with the current policy, providing a
natural curriculum.

\subsection{AlphaGo and AlphaZero}

The most famous application of self-play is the AlphaGo and
AlphaZero\index{AlphaZero|textbf} programme (Silver et al., 2016;
2017). Go\index{Go (game)} had long resisted AI solutions because its
game tree has a branching factor of about $250$ and depths exceeding
$200$, making exhaustive search infeasible.

AlphaGo combined three components: (1) a value network $v_\theta(s)$
trained by regression on self-play outcomes; (2) a policy network
$\pi_\theta(a \mid s)$ trained by supervised learning on human games
and then refined by RL policy gradient; and (3) Monte Carlo Tree
Search (MCTS)\index{MCTS} that uses both networks to guide rollouts.
At each MCTS node, the selection step balances exploration and
exploitation via the PUCT criterion\index{PUCT}:
\begin{equation}
\label{eq:puct}
a^* = \argmax_{a}\!\left[
  Q(s, a) + c_\mathrm{puct}\, \pi_\theta(a \mid s)
  \frac{\sqrt{N(s)}}{1 + N(s, a)}
\right],
\end{equation}
where $Q(s,a)$ is the empirical mean value from past rollouts,
$N(s)$ is the visit count of node $s$, and $N(s,a)$ is the count
of edge $(s,a)$.

AlphaZero\index{AlphaZero} stripped away the supervised learning
bootstrap: starting from \emph{random play}, self-play plus MCTS
alone was sufficient to reach superhuman performance in Go, Chess,
and Shogi within days of training. The success of AlphaZero
demonstrated that given a perfect simulator, self-play can discover
expert strategies entirely from scratch\,---\,a landmark result for
both AI and MARL.

\subsection{Competitive Training Instability and Non-Transitive Strategies}

Naive self-play is not without pathologies. In games with
\textbf{non-transitive strategy
relationships}\index{non-transitive strategies}\,---\,where strategy
$A$ beats $B$, $B$ beats $C$, but $C$ beats $A$ (a rock-paper-scissors
dynamic)\,---\,self-play can cycle indefinitely rather than converging
to a Nash equilibrium. The policy trained against version $k$ of itself
may be strong against $k$ but brittle against version $k-5$, a problem
known as \textbf{forgetting}\index{forgetting!in self-play} or
strategic collapse.

Two remedies are widely used. \textbf{Fictitious
self-play}\index{fictitious self-play} trains the agent against the
\emph{average} of all past opponent policies rather than just the most
recent checkpoint. In two-player zero-sum games with finite action
spaces, fictitious self-play converges to the Nash equilibrium.
\textbf{Population-based training}\index{population-based training}
(PBT; Jaderberg et al., 2017) maintains a diverse league of agents with
different strategies; each agent in the population is trained against
a mixture of others, preventing any single dominant strategy from
eliminating diversity.

% ===========================================================
\section{Communication and Emergent Behaviour}
\label{sec:communication}
% ===========================================================

\subsection{Learned Communication Protocols}

When agents can exchange messages before or during task execution,
the team's effective policy space expands dramatically. But what
should they say? Hard-coding a communication protocol requires
domain expertise and scales poorly. A more principled approach is
to treat messages as \emph{continuous actions} and let communication
protocols emerge through end-to-end training.

\textbf{CommNet}\index{CommNet|textbf} (Sukhbaatar et al., 2016)
is one of the earliest architectures for differentiable multi-agent
communication. At each timestep each agent $i$ broadcasts a
continuous vector $m^i_t \in \R^d$; the aggregated message
$\bar{m}^i_t = \frac{1}{N-1}\sum_{j \neq i} m^j_t$ is concatenated
with the agent's observation and fed to an LSTM\index{LSTM}, whose
output drives both the next action and the next message:
\[
  (a^i_t, m^i_{t+1}) = f_\theta\!\bigl(h^i_t, o^i_t, \bar{m}^i_t\bigr).
\]
Because the communication channel is differentiable, the whole
system is trained end-to-end with standard backpropagation through
time. CommNet demonstrated that agents could spontaneously learn to
coordinate on cooperative tasks (e.g.\ traffic junction routing)
without any supervision on the communication content.

\textbf{TarMAC}\index{TarMAC|textbf} (Das et al., 2019) extended
CommNet by introducing \emph{targeted communication}: each agent
produces a key $k^i$ and a value $v^i$, and the aggregated message
received by agent $j$ is an attention-weighted sum
\[
  m^j = \sum_{i \neq j} \sigma(k^i \cdot q^j)\, v^i,
\]
where $q^j$ is agent $j$'s query vector. The attention mechanism
allows agents to selectively read the messages most relevant to their
current situation, improving both performance and the
interpretability of the learned communication.

\subsection{Emergent Communication}

A striking finding in the emergent-communication literature is that
agents can develop \emph{compositional languages}\index{emergent
language!compositionality}: short discrete symbols whose combinations
encode structured meanings. In referential games\index{referential game}
(e.g.\ one agent sees an object and must describe it; another agent
must identify it from distractors), trained agents frequently develop
vocabularies in which different symbols correspond to different object
attributes, and novel combinations of symbols denote novel objects.
These findings parallel the study of language evolution and raise the
question of whether artificial emergent languages could bootstrap
useful communication skills for real-world agents.

\subsection{Social Dilemmas in Multi-Agent Learning}

The Prisoner's Dilemma of Example~\ref{ex:prisoners} is a one-shot
game. In \emph{repeated} games, cooperation can emerge spontaneously
because agents have the ability to \emph{punish} defectors in future
rounds. Axelrod's famous tournaments showed that the simple
\textsc{Tit-for-Tat} strategy\index{Tit-for-Tat} (cooperate
initially; on each subsequent round, copy the opponent's last action)
outperforms complex strategies in repeated Prisoner's Dilemmas.

In MARL, social dilemmas\index{social dilemma} are formalised as
environments where the Nash equilibrium is Pareto suboptimal: rational
individual behaviour leads to collectively poor outcomes. The
sequential social dilemma (SSD)\index{sequential social dilemma}
framework (Leibo et al., 2017) embeds such dilemmas in spatial
environments. For example, in the \emph{Gathering} game, two agents
collect apples and can fire ``fining beams'' to temporarily disable
each other. Mutual fining wastes time, but each agent is individually
incentivised to fire first. RL agents trained with independent
learning converge to mutual defection; agents trained with prosocial
reward shaping or explicit communication learn to cooperate.

% ===========================================================
\section{Multi-Agent RL for LLMs}
\label{sec:marl-llms}
% ===========================================================

The preceding sections treated agents as neural networks that map
observations to actions in an environment with clear state and reward
signals. Large language models (LLMs)\index{large language model}
introduce a qualitatively new kind of agent: one that operates in the
space of natural language, can reason through intermediate steps, and
can delegate sub-tasks to other LLM instances. MARL concepts apply
naturally, and a growing body of work exploits them to improve
reasoning quality, code generation, and task completion.

\subsection{Debate}

\textbf{AI Debate}\index{debate (MARL)} (Irving et al., 2018) is a
multi-agent protocol for producing trustworthy outputs from a
powerful AI system. Two language model agents argue opposite sides of
a proposition; a third agent (or a human judge) evaluates the
exchange and declares a winner. The key theoretical claim is that,
for a sufficiently capable judge, the optimal strategy for each debater
is to \emph{tell the truth}: a lie can always be exposed by the
opponent, so the best policy in a debate equilibrium is honest argument.

In practice, debate-style protocols are implemented as a two-LLM
pipeline:
\begin{enumerate}[leftmargin=*]
  \item A \textbf{proposer} LLM generates a claim or answer.
  \item A \textbf{critic} LLM, given the same prompt and the
    proposer's output, generates a counter-argument.
  \item A \textbf{judge} LLM (or the same model with a judge prompt)
    scores both positions and selects the better answer.
\end{enumerate}
The judge's score can be used as a reward signal to fine-tune both
the proposer and critic via RL (Section~\ref{sec:rlvr} in
Chapter~\ref{ch:reasoning}). This enables \emph{scalable oversight}:
the system can improve beyond the judge's direct capability,
because it is easier for the judge to evaluate an argument than
to generate the correct answer from scratch.

\subsection{Multi-Agent Code Generation}

Code generation is a natural setting for multi-agent LLMs because
code admits automatic, verifiable rewards (test suites), and because
the task decomposes into sub-tasks (specification, implementation,
testing, documentation) that are easy to assign to specialised agents.

A representative pipeline (Chen et al., 2023; Du et al., 2023) works
as follows:
\begin{enumerate}[leftmargin=*]
  \item A \textbf{planner} agent decomposes the user request into a
    sequence of subtasks.
  \item Multiple \textbf{coder} agents independently generate
    candidate implementations of each subtask.
  \item A \textbf{verifier} agent runs test suites on each candidate
    and returns pass/fail feedback.
  \item A \textbf{ranker} agent, trained with preference RL, selects
    the best-ranked implementation and optionally patches failures.
\end{enumerate}
Empirically, diverse proposals from multiple independent coders
consistently outperform any single coder, even when the individual
coders are identical models: diversity in random seeds, prompts, or
temperature settings introduces sufficient variation for the verifier
to identify genuinely better solutions.

\subsection{LLM Agents Cooperating on Tasks}

Beyond code, multi-agent LLM frameworks have been applied to research
assistance, game playing, and software engineering. \textbf{ChatDev}
(Qian et al., 2023) structures a software team as a set of LLM agents
playing different roles (CEO, CTO, programmer, tester, document writer)
who communicate via a moderated chat channel to complete a software
project end-to-end. Each role is defined by a system prompt and
a role-specific history; the moderation protocol enforces turn-taking
and task focus.

\textbf{AutoGen}\index{AutoGen} (Wu et al., 2023) provides a more
flexible framework: any LLM can be wrapped as an \emph{agent} with a
name, system prompt, and optionally a tool-use interface. Agents
communicate by passing messages in a shared thread; a human-proxy
agent can inject user feedback. The framework supports loops, branching
conversations, and nested agent groups, enabling both cooperative and
adversarial patterns.

\subsection{Agentic Workflows with Specialised Agents}

The most practically influential multi-agent LLM pattern at the time
of writing is the \textbf{orchestrator\,---\,specialist} pattern.
A powerful generalist LLM (the \emph{orchestrator}) receives a complex
user task, decomposes it into subtasks, and dispatches each subtask to
a \emph{specialist} agent optimised for that specific skill:
\begin{description}[leftmargin=!, labelwidth=3.5cm]
  \item[Web search agent] retrieves and summarises live information.
  \item[Code agent] writes and executes code in a sandboxed
    interpreter.
  \item[Memory agent] maintains a persistent key-value store of facts
    extracted from prior turns.
  \item[Critic agent] reviews the orchestrator's plan and flags
    inconsistencies before execution.
\end{description}
Rewards for RL fine-tuning of the orchestrator can come from
end-to-end task completion (e.g.\ does the final answer match a
ground-truth solution?) or from intermediate sub-goals validated by
the specialist agents. This closed loop between orchestration and
specialist execution mirrors the CTDE paradigm: the orchestrator
acts as a centralised planner during training, while each specialist
executes its local policy independently.

The connection to MARL is direct. Each specialist is a policy
$\pi^i$; the orchestrator defines a task distribution over joint
action sequences; and the emergent behaviour of the ensemble is a
trajectory in the space of all possible multi-agent interactions.
Formally training this system with MARL objectives\,---\,rather than
fine-tuning each component in isolation\,---\,is an active research
frontier.

% ===========================================================
\section{Challenges in Multi-Agent RL}
\label{sec:marl-challenges}
% ===========================================================

\subsection{Non-Stationarity}

The most fundamental challenge of MARL is that each agent's
environment includes all other agents, which are themselves
learning. From agent $i$'s perspective, the transition law and
reward function are \emph{non-stationary}\index{non-stationarity!in MARL|textbf}:
the world it experiences today is different from the world it will
experience tomorrow, because its teammates or opponents have updated
their policies. This violates the stationarity assumption that
underpins convergence proofs for Q-learning and other single-agent
algorithms.

Formally, agent $i$ observes an effective transition law
\[
  \tilde{\cT}^i(s' \mid s, a^i)
  = \sum_{\bm{a}^{-i}} \cT(s' \mid s, a^i, \bm{a}^{-i})
    \prod_{j \neq i} \pi^j(a^j \mid s),
\]
which changes whenever any $\pi^j$ ($j \neq i$) changes. Independent
Q-learning\index{independent Q-learning} ignores this non-stationarity
and often fails to converge in practice, especially in competitive
settings. Algorithms such as QMIX and MAPPO partially mitigate
non-stationarity by sharing gradient information across agents, but
convergence guarantees remain weak for the general case.

\subsection{Credit Assignment in Teams}

Even with a shared reward signal, team credit assignment is
hard\index{credit assignment!in MARL}: if the team succeeds, which
agent's action was responsible? Naive approaches attribute all reward
to all agents equally, which provides noisy gradients and discourages
specialisation. The \textbf{difference reward}\index{difference reward}
heuristic assigns agent $i$ the credit
\[
  d^i_t = r_t - r_t^{-i},
\]
where $r_t^{-i}$ is the reward that \emph{would have been} obtained
had agent $i$ acted with some default or counterfactual action,
holding other agents fixed. Computing $r_t^{-i}$ requires either a
learned world model or additional rollouts, making it expensive.
QMIX's monotonicity constraint provides an implicit credit assignment
mechanism: the gradient of $Q_\mathrm{tot}$ with respect to $Q^i$
tells agent $i$ how much its utility function contributes to the team.

\subsection{Scalability}

The joint action space $|\cA| = \prod_i |\cA^i|$ grows exponentially
in the number of agents. For $N = 10$ agents each with $|\cA^i| = 5$
actions, the joint space has $5^{10} \approx 10^7$ elements. Even
representing $Q_\mathrm{tot}(s, \bm{a})$ as a table is infeasible;
the CTDE factorisation of QMIX addresses this by writing
$Q_\mathrm{tot}$ as a function of per-agent utilities, reducing
complexity from exponential to linear in $N$.

Policy-gradient methods scale better because they operate on
marginal policies $\pi^i(a^i \mid \tau^i)$ rather than on the joint
policy. With parameter sharing, the same network services all agents,
so the parameter count does not grow with $N$. Nevertheless,
training time still scales with the number of agents because more
rollout data is needed, and inter-agent dependencies require larger
replay buffers and mini-batches.

\subsection{Reward Shaping for Cooperation}

In mixed-motive games, agents may have individually rational
incentives that conflict with team performance. Reward shaping
\index{reward shaping!in MARL}\,---\,augmenting the true reward
with a potential-based bonus\,---\,can nudge agents toward
cooperative behaviour without invalidating the Nash equilibrium.
However, shaping rewards poorly can introduce spurious equilibria
or make agents indifferent to the true objective.

\textbf{Intrinsic social motivation}\index{intrinsic motivation!social}
is a promising alternative: agents are given an intrinsic reward
for \emph{influencing} the behaviour of others (Social Influence
reward; Jaques et al., 2019) or for achieving causal empowerment
over the environment. These intrinsic rewards encourage information
sharing without requiring hand-crafted communication protocols.

\subsection{Sample Efficiency}

MARL is sample-hungry. Not only must each agent explore its own
action space, but the joint exploration problem is exponentially
larger. Coordinated exploration\,---\,where agents actively
avoid redundant exploration of the same state regions\,---\,reduces
redundancy but requires additional coordination overhead. Offline
MARL\index{offline MARL}, which learns from pre-collected datasets
without online interaction, is an emerging approach that transfers
recent successes of offline single-agent RL
(Chapter~\ref{ch:open}) to the multi-agent setting.

The table below summarises the main algorithmic families,
their settings, and key properties.

\begin{center}
\small
\begin{tabular}{@{}lllll@{}}
\toprule
\textbf{Algorithm} & \textbf{Setting} & \textbf{Paradigm}
  & \textbf{Action space} & \textbf{Key idea} \\
\midrule
IQL & Coop./Comp. & Decentralised & Discrete & Independent Q-learning \\
QMIX & Cooperative & CTDE & Discrete & Monotonic mixer \\
QPLEX & Cooperative & CTDE & Discrete & Duplex attention mixer \\
MAPPO & Cooperative & CTDE & Cont./Disc. & Shared centralised critic \\
MADDPG & Mixed & CTDE & Continuous & Per-agent centralised critic \\
AlphaZero & Competitive & Self-play & Discrete & MCTS + policy/value nets \\
PSRO & Mixed & Population & Discrete & Nash resp.\ oracle + mixture \\
CommNet & Cooperative & Comm.\ channel & Continuous & Mean-field messages \\
TarMAC & Cooperative & Comm.\ channel & Continuous & Targeted attention msgs \\
\bottomrule
\end{tabular}
\end{center}

% ===========================================================
\section{Summary}
\label{sec:marl-summary}
% ===========================================================

Multi-agent reinforcement learning extends the single-agent MDP
framework to $N$ interacting agents, each with its own policy,
observation, and reward. The central solution concept is Nash
equilibrium: a joint policy from which no agent can profitably
deviate unilaterally. Finding Nash equilibria is hard in general,
but structure in the game\,---\,cooperation, zero-sum competition,
factored dynamics\,---\,enables tractable algorithms.

For cooperative tasks, the CTDE paradigm achieves the best of
both worlds: centralised gradient information at training time,
and fully decentralised execution. QMIX exploits monotonic
factorisation to make greedy per-agent action selection globally
optimal. MAPPO applies PPO with a centralised value baseline and
parameter sharing, scaling to large teams with continuous actions.

In competitive settings, self-play provides an automatic curriculum.
AlphaZero shows that starting from random play and iterating
self-play with MCTS is sufficient for superhuman performance in
perfect-information games. Non-transitivity and strategic forgetting
motivate population-based methods that maintain policy diversity.

Communication opens an additional channel for cooperation.
CommNet and TarMAC demonstrate that end-to-end training can
produce effective messaging protocols without any supervision
on the message content. In social dilemmas, where the Nash
equilibrium is Pareto suboptimal, repeated interaction and
prosocial rewards can sustain cooperative behaviour.

LLMs are the latest and perhaps most impactful arena for MARL
ideas. Debate, multi-agent code generation, orchestrator-specialist
pipelines, and frameworks like AutoGen and ChatDev all instantiate
multi-agent dynamics in language. Applying formal MARL training
objectives to these systems\,---\,rather than fine-tuning components
in isolation\,---\,is one of the most exciting open directions at
the frontier of RL research.

The key challenges\,---\,non-stationarity, credit assignment,
scalability, reward shaping, and sample efficiency\,---\,are
interrelated and none is fully solved. Progress on each will likely
require new theory as well as new engineering, making MARL one of
the most active and consequential subfields of machine learning.

% ===========================================================
\section*{Exercises}
\addcontentsline{toc}{section}{Exercises}
% ===========================================================

\begin{exercise}
\label{ex:markov-game-reduction}
Show that a Markov game with $N=1$ reduces to a standard MDP.
Then show that if all agents share the same reward $R^1 = \cdots =
R^N$ and all agents act identically (i.e.\ $\pi^1 = \cdots = \pi^N$),
the joint return Equation~\eqref{eq:cooperative-return} equals the
single-agent return of Chapter~\ref{ch:mdp}. What additional
conditions are needed for the shared-reward policy-gradient to equal
the single-agent policy-gradient?
\end{exercise}

\begin{exercise}
\label{ex:nash-prisoners}
Consider the Prisoner's Dilemma of Example~\ref{ex:prisoners}.
\begin{enumerate}[label=(\alph*)]
  \item Verify that $(D, D)$ is a Nash equilibrium by checking the
    definition: no player can improve their payoff by unilaterally
    switching to $C$.
  \item Show that $(C, C)$ is \emph{not} a Nash equilibrium.
  \item Now suppose the game is played for $T = 100$ rounds and
    agents maximise total undiscounted reward. Is it a Nash
    equilibrium for both agents to always cooperate? Justify.
  \item What changes if the horizon is infinite with discount
    $\gamma = 0.9$?
\end{enumerate}
\end{exercise}

\begin{exercise}
\label{ex:igm-proof}
The IGM property states that
$\argmax_{\bm{a}} Q_\mathrm{tot}(s, \bm{a}) = \bigl(\argmax_{a^1}
Q^1(\tau^1, a^1),\, \ldots,\, \argmax_{a^N} Q^N(\tau^N,
a^N)\bigr)$.
\begin{enumerate}[label=(\alph*)]
  \item Suppose $Q_\mathrm{tot}(s, \bm{a}) = \sum_i Q^i(\tau^i,
    a^i)$. Prove that the IGM property holds.
  \item Does the IGM property hold if
    $Q_\mathrm{tot}(s, \bm{a}) = \prod_i Q^i(\tau^i, a^i)$ and
    all $Q^i > 0$? Prove or give a counterexample.
  \item Explain in one paragraph why QMIX's non-negative mixing
    weights (from its hypernetwork) are sufficient but not necessary
    for IGM.
\end{enumerate}
\end{exercise}

\begin{exercise}
\label{ex:self-play-rps}
Consider a three-strategy game with pure strategies Rock (R), Paper
(P), Scissors (S) and payoff matrix
\[
M = \begin{pmatrix} 0 & -1 & 1 \\ 1 & 0 & -1 \\ -1 & 1 & 0
\end{pmatrix}
\]
where $M_{ij}$ is the payoff to the row player when they play
strategy $i$ against column player strategy $j$.
\begin{enumerate}[label=(\alph*)]
  \item Find the unique Nash equilibrium of this game.
  \item Simulate $1{,}000$ steps of gradient ascent on $\pi^1$ while
    $\pi^2$ is fixed at the Nash equilibrium. Does $\pi^1$ converge?
  \item Now simulate $1{,}000$ steps of simultaneous gradient ascent
    for both players. Describe the trajectory in the space of mixed
    strategies. Does it converge to the Nash equilibrium?
  \item Explain intuitively why this game causes naive self-play to
    cycle.
\end{enumerate}
\end{exercise}

\begin{exercise}
\label{ex:credit-assignment}
A cooperative team of $N = 4$ robots earns a shared reward of $r =
10$ when all four reach a goal simultaneously, and $r = 0$ otherwise.
After a successful episode, independent Q-learning assigns equal
credit $10/4 = 2.5$ to each agent.
\begin{enumerate}[label=(\alph*)]
  \item Suppose Robot 3 had a $0.5$ chance of reaching the goal
    regardless of its action. Describe how equal credit assignment
    distorts the learning signal for Robots 1, 2, and 4.
  \item Define a difference reward for Robot $i$ using a counterfactual
    where Robot $i$ stays put while others follow the observed policy.
    Compute this reward for Robot 3.
  \item What is the advantage of the difference reward over equal
    credit? What is its main computational drawback?
\end{enumerate}
\end{exercise}

\begin{exercise}
\label{ex:ctde-communication}
Design a CTDE algorithm for a team of $N$ drones surveying a
search-and-rescue area. Each drone observes a $5 \times 5$ local
grid; the global map is $50 \times 50$. Agents cannot communicate
at test time.
\begin{enumerate}[label=(\alph*)]
  \item Specify the observation space $\mathcal{O}^i$, action space
    $\mathcal{A}^i$, and a suitable team reward function $R$.
  \item Propose a centralised critic architecture that uses the
    full $50 \times 50$ map during training. What features would
    you concatenate alongside the global state?
  \item Argue whether parameter sharing is appropriate here, and
    how you would condition each agent's actor on its unique
    identifier.
  \item Identify two ways in which training instability might arise
    and suggest mitigations.
\end{enumerate}
\end{exercise}

\begin{exercise}
\label{ex:debate-reward}
Implement a simplified debate protocol for a closed-book
question-answering task.
\begin{enumerate}[label=(\alph*)]
  \item Define the state, action, and reward for each of the
    proposer, critic, and judge agents in a single-round debate.
  \item Let the judge assign a binary reward $r \in \{0, 1\}$
    (correct / incorrect). Write the REINFORCE policy gradient
    update for the proposer, assuming the critic's policy is
    fixed. Is the variance of this estimator likely to be high?
    Why?
  \item Propose a baseline for the proposer that reduces variance
    without introducing bias. Describe concretely how you would
    estimate it.
  \item Under what conditions does the debate equilibrium guarantee
    truthful behaviour from the proposer? State the necessary
    assumptions on the judge's capability.
\end{enumerate}
\end{exercise}

\begin{exercise}
\label{ex:emergent-language}
In a referential game, a speaker agent observes an image drawn from
a set of $K$ objects and must produce a discrete message
$m \in \{1, \ldots, V\}^L$ (vocabulary size $V$, message length $L$).
A listener observes $m$ and must identify the correct object from $K$
distractors.
\begin{enumerate}[label=(\alph*)]
  \item What is the maximum number of distinct objects that can be
    encoded by a language with vocabulary $V$ and length $L$?
  \item Suppose both agents are trained end-to-end with REINFORCE
    and the reward is $+1$ for correct identification, $0$ otherwise.
    Write the joint policy gradient update.
  \item Describe one metric for measuring the compositionality of
    the emergent language (e.g.\ topographic similarity). Define it
    formally.
  \item Why might the emergent language fail to be compositional
    even when compositionality would be useful?
\end{enumerate}
\end{exercise}
